%% macros_kit.tex -- distilled ICML math + text macro kit for the icml-paper skill.
%%
%% This is an ORIGINAL curation of conventional LaTeX math-writing patterns,
%% distilled from two exemplar sources and stripped of each paper's project-
%% specific notation. Nothing here is a verbatim copy of either paper's kit.
%%   - Paper A "The Flexibility Trap" (config.tex): text shortcuts + run-in headers.
%%   - Paper B "High-accuracy sampling ..." (macros.tex, arxiv_style.tex): paired
%%     delimiters, calligraphic alphabet, expectation/divergence operators, and the
%%     shared-counter theorem environments.
%%
%% Load these packages in the preamble BEFORE %% macros_kit.tex -- distilled ICML math + text macro kit for the icml-paper skill.
%%
%% This is an ORIGINAL curation of conventional LaTeX math-writing patterns,
%% distilled from two exemplar sources and stripped of each paper's project-
%% specific notation. Nothing here is a verbatim copy of either paper's kit.
%%   - Paper A "The Flexibility Trap" (config.tex): text shortcuts + run-in headers.
%%   - Paper B "High-accuracy sampling ..." (macros.tex, arxiv_style.tex): paired
%%     delimiters, calligraphic alphabet, expectation/divergence operators, and the
%%     shared-counter theorem environments.
%%
%% Load these packages in the preamble BEFORE %% macros_kit.tex -- distilled ICML math + text macro kit for the icml-paper skill.
%%
%% This is an ORIGINAL curation of conventional LaTeX math-writing patterns,
%% distilled from two exemplar sources and stripped of each paper's project-
%% specific notation. Nothing here is a verbatim copy of either paper's kit.
%%   - Paper A "The Flexibility Trap" (config.tex): text shortcuts + run-in headers.
%%   - Paper B "High-accuracy sampling ..." (macros.tex, arxiv_style.tex): paired
%%     delimiters, calligraphic alphabet, expectation/divergence operators, and the
%%     shared-counter theorem environments.
%%
%% Load these packages in the preamble BEFORE %% macros_kit.tex -- distilled ICML math + text macro kit for the icml-paper skill.
%%
%% This is an ORIGINAL curation of conventional LaTeX math-writing patterns,
%% distilled from two exemplar sources and stripped of each paper's project-
%% specific notation. Nothing here is a verbatim copy of either paper's kit.
%%   - Paper A "The Flexibility Trap" (config.tex): text shortcuts + run-in headers.
%%   - Paper B "High-accuracy sampling ..." (macros.tex, arxiv_style.tex): paired
%%     delimiters, calligraphic alphabet, expectation/divergence operators, and the
%%     shared-counter theorem environments.
%%
%% Load these packages in the preamble BEFORE \input{macros_kit}:
%%   amsmath, amssymb, mathtools, amsthm, xspace
%% Optional: bbm (for a blackboard indicator; then switch \ind to \mathbbm{1}).

% ---- Paired delimiters (use the starred form for auto-sizing, e.g. \norm*{x}) ----
\DeclarePairedDelimiter{\abs}{\lvert}{\rvert}
\DeclarePairedDelimiter{\norm}{\lVert}{\rVert}
\DeclarePairedDelimiter{\prn}{(}{)}
\DeclarePairedDelimiter{\brk}{[}{]}
\DeclarePairedDelimiter{\crl}{\{}{\}}
\DeclarePairedDelimiter{\ceil}{\lceil}{\rceil}
\DeclarePairedDelimiter{\floor}{\lfloor}{\rfloor}
\DeclarePairedDelimiter{\inner}{\langle}{\rangle}
\newcommand{\set}[1]{\crl*{#1}}

% ---- Blackboard number sets (explicit safe names; no risky single-letter clobber) ----
\newcommand{\R}{\mathbb{R}}
\newcommand{\N}{\mathbb{N}}
\newcommand{\Z}{\mathbb{Z}}
\newcommand{\Q}{\mathbb{Q}}
\newcommand{\C}{\mathbb{C}}

% ---- Calligraphic alphabet \cA .. \cZ (the one auto-loop worth keeping) ----
\newcommand{\defcalX}[1]{\expandafter\newcommand\csname c#1\endcsname{\mathcal{#1}}}
\defcalX{A}\defcalX{B}\defcalX{C}\defcalX{D}\defcalX{E}\defcalX{F}\defcalX{G}
\defcalX{H}\defcalX{I}\defcalX{J}\defcalX{K}\defcalX{L}\defcalX{M}\defcalX{N}
\defcalX{O}\defcalX{P}\defcalX{Q}\defcalX{R}\defcalX{S}\defcalX{T}\defcalX{U}
\defcalX{V}\defcalX{W}\defcalX{X}\defcalX{Y}\defcalX{Z}
\let\defcalX\undefined

% ---- Operators, expectation, probability ----
\DeclareMathOperator*{\argmin}{arg\,min}
\DeclareMathOperator*{\argmax}{arg\,max}
\DeclareMathOperator{\Tr}{Tr}
\DeclareMathOperator{\Var}{Var}
\DeclareMathOperator{\Cov}{Cov}
\DeclareMathOperator{\supp}{supp}
\newcommand{\E}{\mathbb{E}}
\newcommand{\Prob}{\mathbb{P}}          % \Prob, not \P (\P is the text pilcrow)
\newcommand{\ind}{\mathbf{1}}           % blackboard variant: load bbm, use \mathbbm{1}
\newcommand{\eps}{\varepsilon}
\newcommand{\T}{^{\top}}
\newcommand{\defeq}{\vcentcolon=}
\newcommand{\eqdef}{=\vcentcolon}

% ---- Hat / tilde / bar wrappers ----
\newcommand{\wt}[1]{\widetilde{#1}}
\newcommand{\wh}[1]{\widehat{#1}}
\newcommand{\wb}[1]{\overline{#1}}

% ---- Divergences and distances (auto-sized via \prn*) ----
\newcommand{\Dkl}[2]{D_{\mathrm{KL}}\prn*{#1 \,\|\, #2}}
\newcommand{\Dtv}[2]{D_{\mathrm{TV}}\prn*{#1, #2}}
\newcommand{\Dbl}[2]{D_{\mathrm{BL}}\prn*{#1, #2}}
\newcommand{\Dhel}[2]{D_{\mathrm{H}}\prn*{#1, #2}}
\newcommand{\Dchi}[2]{D_{\chi^2}\prn*{#1 \,\|\, #2}}

% ---- Text shortcuts (Paper A idiom) ----
\newcommand{\ie}{i.e.,\xspace}
\newcommand{\eg}{e.g.,\xspace}
\newcommand{\cf}{cf.\xspace}
\newcommand{\etal}{et~al.\xspace}
\newcommand{\etc}{etc.\xspace}
\newcommand{\wrt}{w.r.t.\xspace}
% Run-in bold header (Paper A style). Named \phead so it does not clobber the
% ICML \paragraph. Usage: \phead{Contributions.} text continues on the same line.
\newcommand{\phead}[1]{\vspace{2pt}\noindent\textbf{#1}\hspace{0.5em}}

% ---- Theorem-like environments (Paper B pattern: one shared counter, per section) ----
\theoremstyle{plain}
\newtheorem{theorem}{Theorem}[section]
\newtheorem{proposition}[theorem]{Proposition}
\newtheorem{lemma}[theorem]{Lemma}
\newtheorem{corollary}[theorem]{Corollary}
\theoremstyle{definition}
\newtheorem{definition}[theorem]{Definition}
\newtheorem{assumption}[theorem]{Assumption}
\theoremstyle{remark}
\newtheorem{remark}[theorem]{Remark}
\newtheorem{example}[theorem]{Example}
:
%%   amsmath, amssymb, mathtools, amsthm, xspace
%% Optional: bbm (for a blackboard indicator; then switch \ind to \mathbbm{1}).

% ---- Paired delimiters (use the starred form for auto-sizing, e.g. \norm*{x}) ----
\DeclarePairedDelimiter{\abs}{\lvert}{\rvert}
\DeclarePairedDelimiter{\norm}{\lVert}{\rVert}
\DeclarePairedDelimiter{\prn}{(}{)}
\DeclarePairedDelimiter{\brk}{[}{]}
\DeclarePairedDelimiter{\crl}{\{}{\}}
\DeclarePairedDelimiter{\ceil}{\lceil}{\rceil}
\DeclarePairedDelimiter{\floor}{\lfloor}{\rfloor}
\DeclarePairedDelimiter{\inner}{\langle}{\rangle}
\newcommand{\set}[1]{\crl*{#1}}

% ---- Blackboard number sets (explicit safe names; no risky single-letter clobber) ----
\newcommand{\R}{\mathbb{R}}
\newcommand{\N}{\mathbb{N}}
\newcommand{\Z}{\mathbb{Z}}
\newcommand{\Q}{\mathbb{Q}}
\newcommand{\C}{\mathbb{C}}

% ---- Calligraphic alphabet \cA .. \cZ (the one auto-loop worth keeping) ----
\newcommand{\defcalX}[1]{\expandafter\newcommand\csname c#1\endcsname{\mathcal{#1}}}
\defcalX{A}\defcalX{B}\defcalX{C}\defcalX{D}\defcalX{E}\defcalX{F}\defcalX{G}
\defcalX{H}\defcalX{I}\defcalX{J}\defcalX{K}\defcalX{L}\defcalX{M}\defcalX{N}
\defcalX{O}\defcalX{P}\defcalX{Q}\defcalX{R}\defcalX{S}\defcalX{T}\defcalX{U}
\defcalX{V}\defcalX{W}\defcalX{X}\defcalX{Y}\defcalX{Z}
\let\defcalX\undefined

% ---- Operators, expectation, probability ----
\DeclareMathOperator*{\argmin}{arg\,min}
\DeclareMathOperator*{\argmax}{arg\,max}
\DeclareMathOperator{\Tr}{Tr}
\DeclareMathOperator{\Var}{Var}
\DeclareMathOperator{\Cov}{Cov}
\DeclareMathOperator{\supp}{supp}
\newcommand{\E}{\mathbb{E}}
\newcommand{\Prob}{\mathbb{P}}          % \Prob, not \P (\P is the text pilcrow)
\newcommand{\ind}{\mathbf{1}}           % blackboard variant: load bbm, use \mathbbm{1}
\newcommand{\eps}{\varepsilon}
\newcommand{\T}{^{\top}}
\newcommand{\defeq}{\vcentcolon=}
\newcommand{\eqdef}{=\vcentcolon}

% ---- Hat / tilde / bar wrappers ----
\newcommand{\wt}[1]{\widetilde{#1}}
\newcommand{\wh}[1]{\widehat{#1}}
\newcommand{\wb}[1]{\overline{#1}}

% ---- Divergences and distances (auto-sized via \prn*) ----
\newcommand{\Dkl}[2]{D_{\mathrm{KL}}\prn*{#1 \,\|\, #2}}
\newcommand{\Dtv}[2]{D_{\mathrm{TV}}\prn*{#1, #2}}
\newcommand{\Dbl}[2]{D_{\mathrm{BL}}\prn*{#1, #2}}
\newcommand{\Dhel}[2]{D_{\mathrm{H}}\prn*{#1, #2}}
\newcommand{\Dchi}[2]{D_{\chi^2}\prn*{#1 \,\|\, #2}}

% ---- Text shortcuts (Paper A idiom) ----
\newcommand{\ie}{i.e.,\xspace}
\newcommand{\eg}{e.g.,\xspace}
\newcommand{\cf}{cf.\xspace}
\newcommand{\etal}{et~al.\xspace}
\newcommand{\etc}{etc.\xspace}
\newcommand{\wrt}{w.r.t.\xspace}
% Run-in bold header (Paper A style). Named \phead so it does not clobber the
% ICML \paragraph. Usage: \phead{Contributions.} text continues on the same line.
\newcommand{\phead}[1]{\vspace{2pt}\noindent\textbf{#1}\hspace{0.5em}}

% ---- Theorem-like environments (Paper B pattern: one shared counter, per section) ----
\theoremstyle{plain}
\newtheorem{theorem}{Theorem}[section]
\newtheorem{proposition}[theorem]{Proposition}
\newtheorem{lemma}[theorem]{Lemma}
\newtheorem{corollary}[theorem]{Corollary}
\theoremstyle{definition}
\newtheorem{definition}[theorem]{Definition}
\newtheorem{assumption}[theorem]{Assumption}
\theoremstyle{remark}
\newtheorem{remark}[theorem]{Remark}
\newtheorem{example}[theorem]{Example}
:
%%   amsmath, amssymb, mathtools, amsthm, xspace
%% Optional: bbm (for a blackboard indicator; then switch \ind to \mathbbm{1}).

% ---- Paired delimiters (use the starred form for auto-sizing, e.g. \norm*{x}) ----
\DeclarePairedDelimiter{\abs}{\lvert}{\rvert}
\DeclarePairedDelimiter{\norm}{\lVert}{\rVert}
\DeclarePairedDelimiter{\prn}{(}{)}
\DeclarePairedDelimiter{\brk}{[}{]}
\DeclarePairedDelimiter{\crl}{\{}{\}}
\DeclarePairedDelimiter{\ceil}{\lceil}{\rceil}
\DeclarePairedDelimiter{\floor}{\lfloor}{\rfloor}
\DeclarePairedDelimiter{\inner}{\langle}{\rangle}
\newcommand{\set}[1]{\crl*{#1}}

% ---- Blackboard number sets (explicit safe names; no risky single-letter clobber) ----
\newcommand{\R}{\mathbb{R}}
\newcommand{\N}{\mathbb{N}}
\newcommand{\Z}{\mathbb{Z}}
\newcommand{\Q}{\mathbb{Q}}
\newcommand{\C}{\mathbb{C}}

% ---- Calligraphic alphabet \cA .. \cZ (the one auto-loop worth keeping) ----
\newcommand{\defcalX}[1]{\expandafter\newcommand\csname c#1\endcsname{\mathcal{#1}}}
\defcalX{A}\defcalX{B}\defcalX{C}\defcalX{D}\defcalX{E}\defcalX{F}\defcalX{G}
\defcalX{H}\defcalX{I}\defcalX{J}\defcalX{K}\defcalX{L}\defcalX{M}\defcalX{N}
\defcalX{O}\defcalX{P}\defcalX{Q}\defcalX{R}\defcalX{S}\defcalX{T}\defcalX{U}
\defcalX{V}\defcalX{W}\defcalX{X}\defcalX{Y}\defcalX{Z}
\let\defcalX\undefined

% ---- Operators, expectation, probability ----
\DeclareMathOperator*{\argmin}{arg\,min}
\DeclareMathOperator*{\argmax}{arg\,max}
\DeclareMathOperator{\Tr}{Tr}
\DeclareMathOperator{\Var}{Var}
\DeclareMathOperator{\Cov}{Cov}
\DeclareMathOperator{\supp}{supp}
\newcommand{\E}{\mathbb{E}}
\newcommand{\Prob}{\mathbb{P}}          % \Prob, not \P (\P is the text pilcrow)
\newcommand{\ind}{\mathbf{1}}           % blackboard variant: load bbm, use \mathbbm{1}
\newcommand{\eps}{\varepsilon}
\newcommand{\T}{^{\top}}
\newcommand{\defeq}{\vcentcolon=}
\newcommand{\eqdef}{=\vcentcolon}

% ---- Hat / tilde / bar wrappers ----
\newcommand{\wt}[1]{\widetilde{#1}}
\newcommand{\wh}[1]{\widehat{#1}}
\newcommand{\wb}[1]{\overline{#1}}

% ---- Divergences and distances (auto-sized via \prn*) ----
\newcommand{\Dkl}[2]{D_{\mathrm{KL}}\prn*{#1 \,\|\, #2}}
\newcommand{\Dtv}[2]{D_{\mathrm{TV}}\prn*{#1, #2}}
\newcommand{\Dbl}[2]{D_{\mathrm{BL}}\prn*{#1, #2}}
\newcommand{\Dhel}[2]{D_{\mathrm{H}}\prn*{#1, #2}}
\newcommand{\Dchi}[2]{D_{\chi^2}\prn*{#1 \,\|\, #2}}

% ---- Text shortcuts (Paper A idiom) ----
\newcommand{\ie}{i.e.,\xspace}
\newcommand{\eg}{e.g.,\xspace}
\newcommand{\cf}{cf.\xspace}
\newcommand{\etal}{et~al.\xspace}
\newcommand{\etc}{etc.\xspace}
\newcommand{\wrt}{w.r.t.\xspace}
% Run-in bold header (Paper A style). Named \phead so it does not clobber the
% ICML \paragraph. Usage: \phead{Contributions.} text continues on the same line.
\newcommand{\phead}[1]{\vspace{2pt}\noindent\textbf{#1}\hspace{0.5em}}

% ---- Theorem-like environments (Paper B pattern: one shared counter, per section) ----
\theoremstyle{plain}
\newtheorem{theorem}{Theorem}[section]
\newtheorem{proposition}[theorem]{Proposition}
\newtheorem{lemma}[theorem]{Lemma}
\newtheorem{corollary}[theorem]{Corollary}
\theoremstyle{definition}
\newtheorem{definition}[theorem]{Definition}
\newtheorem{assumption}[theorem]{Assumption}
\theoremstyle{remark}
\newtheorem{remark}[theorem]{Remark}
\newtheorem{example}[theorem]{Example}
:
%%   amsmath, amssymb, mathtools, amsthm, xspace
%% Optional: bbm (for a blackboard indicator; then switch \ind to \mathbbm{1}).

% ---- Paired delimiters (use the starred form for auto-sizing, e.g. \norm*{x}) ----
\DeclarePairedDelimiter{\abs}{\lvert}{\rvert}
\DeclarePairedDelimiter{\norm}{\lVert}{\rVert}
\DeclarePairedDelimiter{\prn}{(}{)}
\DeclarePairedDelimiter{\brk}{[}{]}
\DeclarePairedDelimiter{\crl}{\{}{\}}
\DeclarePairedDelimiter{\ceil}{\lceil}{\rceil}
\DeclarePairedDelimiter{\floor}{\lfloor}{\rfloor}
\DeclarePairedDelimiter{\inner}{\langle}{\rangle}
\newcommand{\set}[1]{\crl*{#1}}

% ---- Blackboard number sets (explicit safe names; no risky single-letter clobber) ----
\newcommand{\R}{\mathbb{R}}
\newcommand{\N}{\mathbb{N}}
\newcommand{\Z}{\mathbb{Z}}
\newcommand{\Q}{\mathbb{Q}}
\newcommand{\C}{\mathbb{C}}

% ---- Calligraphic alphabet \cA .. \cZ (the one auto-loop worth keeping) ----
\newcommand{\defcalX}[1]{\expandafter\newcommand\csname c#1\endcsname{\mathcal{#1}}}
\defcalX{A}\defcalX{B}\defcalX{C}\defcalX{D}\defcalX{E}\defcalX{F}\defcalX{G}
\defcalX{H}\defcalX{I}\defcalX{J}\defcalX{K}\defcalX{L}\defcalX{M}\defcalX{N}
\defcalX{O}\defcalX{P}\defcalX{Q}\defcalX{R}\defcalX{S}\defcalX{T}\defcalX{U}
\defcalX{V}\defcalX{W}\defcalX{X}\defcalX{Y}\defcalX{Z}
\let\defcalX\undefined

% ---- Operators, expectation, probability ----
\DeclareMathOperator*{\argmin}{arg\,min}
\DeclareMathOperator*{\argmax}{arg\,max}
\DeclareMathOperator{\Tr}{Tr}
\DeclareMathOperator{\Var}{Var}
\DeclareMathOperator{\Cov}{Cov}
\DeclareMathOperator{\supp}{supp}
\newcommand{\E}{\mathbb{E}}
\newcommand{\Prob}{\mathbb{P}}          % \Prob, not \P (\P is the text pilcrow)
\newcommand{\ind}{\mathbf{1}}           % blackboard variant: load bbm, use \mathbbm{1}
\newcommand{\eps}{\varepsilon}
\newcommand{\T}{^{\top}}
\newcommand{\defeq}{\vcentcolon=}
\newcommand{\eqdef}{=\vcentcolon}

% ---- Hat / tilde / bar wrappers ----
\newcommand{\wt}[1]{\widetilde{#1}}
\newcommand{\wh}[1]{\widehat{#1}}
\newcommand{\wb}[1]{\overline{#1}}

% ---- Divergences and distances (auto-sized via \prn*) ----
\newcommand{\Dkl}[2]{D_{\mathrm{KL}}\prn*{#1 \,\|\, #2}}
\newcommand{\Dtv}[2]{D_{\mathrm{TV}}\prn*{#1, #2}}
\newcommand{\Dbl}[2]{D_{\mathrm{BL}}\prn*{#1, #2}}
\newcommand{\Dhel}[2]{D_{\mathrm{H}}\prn*{#1, #2}}
\newcommand{\Dchi}[2]{D_{\chi^2}\prn*{#1 \,\|\, #2}}

% ---- Text shortcuts (Paper A idiom) ----
\newcommand{\ie}{i.e.,\xspace}
\newcommand{\eg}{e.g.,\xspace}
\newcommand{\cf}{cf.\xspace}
\newcommand{\etal}{et~al.\xspace}
\newcommand{\etc}{etc.\xspace}
\newcommand{\wrt}{w.r.t.\xspace}
% Run-in bold header (Paper A style). Named \phead so it does not clobber the
% ICML \paragraph. Usage: \phead{Contributions.} text continues on the same line.
\newcommand{\phead}[1]{\vspace{2pt}\noindent\textbf{#1}\hspace{0.5em}}

% ---- Theorem-like environments (Paper B pattern: one shared counter, per section) ----
\theoremstyle{plain}
\newtheorem{theorem}{Theorem}[section]
\newtheorem{proposition}[theorem]{Proposition}
\newtheorem{lemma}[theorem]{Lemma}
\newtheorem{corollary}[theorem]{Corollary}
\theoremstyle{definition}
\newtheorem{definition}[theorem]{Definition}
\newtheorem{assumption}[theorem]{Assumption}
\theoremstyle{remark}
\newtheorem{remark}[theorem]{Remark}
\newtheorem{example}[theorem]{Example}
